\chapter{Fluida yang Ingin Kita Tiru}
\label{ch:fluida}

Kuliah metode numerik dalam mekanika fluida saya ikuti pada musim dingin 2025/26, diberikan oleh
Stefan Pirker, yang kemudian juga membimbing tugas akhir saya. Gaya kuliahnya khas. Setiap topik
dibuka dengan pertanyaan seperti ``aliran apa yang terlintas di kepala Anda?'', dan beberapa slide
meminta ChatGPT menulis puisi tentang neraca massa atau pidato perpisahan untuk persamaan
Navier--Stokes, sebelum kembali ke persamaan yang sesungguhnya. Bab ini dan bab berikutnya merangkum
kuliah itu. Bagian III secara keseluruhan memakai fluida sebagai contoh utama, karena di sanalah
saya paling banyak bekerja, tetapi gagasan-gagasannya berlaku untuk banyak sistem fisika lain.

Sebelum membicarakan bagaimana machine learning bisa membantu simulasi, kita perlu tahu apa yang
disimulasikan. Bab ini membahas persamaan yang mengatur aliran fluida dan panas, serta bilangan tak
berdimensi yang menentukan wataknya.

\section{Wajah-wajah aliran}

Aliran fluida punya banyak wajah. Madu yang dituang mengalir pelan dan teratur. Asap rokok naik lurus
sebentar, lalu tiba-tiba pecah menjadi pusaran yang kacau. Awan di belakang pulau kecil di tengah
samudra membentuk deretan pusaran yang bergantian ke kiri dan ke kanan, dan foto satelit pulau
Alexander Selkirk di Pasifik selatan menjadi contoh yang sering dipakai. Jembatan Tacoma Narrows
runtuh pada 1940 karena angin yang berinteraksi dengan getaran strukturnya. Semua ini diatur oleh
persamaan yang sama. Yang membedakannya adalah perbandingan antara gaya-gaya yang bekerja.

Bilangan Reynolds menyatakan perbandingan antara inersia dan gesekan kental,
\begin{equation}
\mathrm{Re}=\frac{UL}{\nu},
\end{equation}
dengan $U$ kecepatan khas, $L$ panjang khas, dan $\nu$ viskositas kinematik. Untuk Re di bawah
sekitar sepuluh, aliran merayap dan gesekan mendominasi. Pada Re menengah, aliran laminar tetapi
bisa melepaskan pusaran secara periodik. Pada Re tinggi, aliran menjadi turbulen. Air dengan
viskositas kinematik sekitar $10^{-6}\,\mathrm{m^2/s}$ yang mengalir satu meter per detik di pipa
berdiameter lima sentimeter punya Re sekitar lima puluh ribu, jauh di daerah turbulen.

Untuk aliran yang digerakkan perbedaan temperatur, bilangan lain menjadi penting. Bilangan Prandtl
membandingkan difusi momentum dengan difusi panas, bilangan Grashof atau Rayleigh mengukur kekuatan
gaya apung dibanding gesekan, dan bilangan Péclet membandingkan perpindahan oleh aliran dengan
perpindahan oleh difusi. Dua aliran dengan bilangan-bilangan tak berdimensi yang sama berperilaku
sama, sekecil atau sebesar apa pun ukurannya. Prinsip kesamaan ini memungkinkan percobaan dengan
model kecil di laboratorium, dan juga menentukan seberapa jauh sebuah model pengganti bisa
diharapkan berlaku.

\section{Hukum kekekalan}

Semua persamaan fluida adalah neraca. Untuk sebuah volume kecil di ruang, perubahan jumlah suatu
besaran di dalamnya sama dengan apa yang mengalir masuk dikurangi apa yang mengalir keluar,
ditambah apa yang diproduksi di dalamnya. Ungkapan matematisnya adalah persamaan transport umum
untuk besaran $\phi$ per satuan massa,
\begin{equation}
\frac{\partial(\rho\phi)}{\partial t}+\nabla\cdot(\rho\vu\phi)=\nabla\cdot(\Gamma\nabla\phi)+S_\phi,
\label{eq:transport}
\end{equation}
dengan suku perubahan terhadap waktu, suku konveksi oleh aliran, suku difusi, dan suku sumber.
Hampir semua persamaan di bab ini adalah kasus khusus bentuk ini.

Kekekalan massa mengambil $\phi=1$, tanpa difusi dan sumber:
\begin{equation}
\frac{\partial\rho}{\partial t}+\nabla\cdot(\rho\vu)=0.
\end{equation}
Untuk fluida tak termampatkan dengan massa jenis tetap, persamaan ini menjadi syarat
$\nabla\cdot\vu=0$: apa yang masuk ke sebuah volume harus keluar lagi.

Kekekalan momentum adalah hukum Newton untuk sebuah paket fluida. Percepatan paket itu dikalikan
massanya sama dengan gaya tekanan, gaya gesekan kental, dan gaya badan seperti gravitasi. Untuk
fluida Newtonian tak termampatkan dengan viskositas tetap, kita mendapat persamaan Navier--Stokes,
\begin{equation}
\rho\Big(\frac{\partial\vu}{\partial t}+(\vu\cdot\nabla)\vu\Big)=-\nabla p+\mu\nabla^2\vu+\rho\bm g.
\label{eq:ns}
\end{equation}
Suku $(\vu\cdot\nabla)\vu$ adalah sumber semua kesulitan. Ia nonlinear, karena kecepatan
mengangkut dirinya sendiri, dan dialah yang melahirkan turbulensi. Tekanan tidak punya persamaan
evolusinya sendiri dalam aliran tak termampatkan. Ia muncul sebagai pengali yang menjaga agar
$\nabla\cdot\vu=0$ terpenuhi di setiap saat, dan kopling inilah yang membuat solver numerik rumit
(\cref{ch:numerik}).

Kekekalan energi untuk aliran dengan perpindahan panas, dalam bentuk temperatur, berbunyi
\begin{equation}
\rho c_p\Big(\frac{\partial T}{\partial t}+\vu\cdot\nabla T\Big)=\nabla\cdot(\lambda\nabla T)+S_T,
\label{eq:energi}
\end{equation}
dengan $c_p$ kapasitas panas, $\lambda$ konduktivitas termal, dan $S_T$ sumber panas seperti reaksi
kimia atau radiasi. Kekekalan spesies untuk fraksi massa $c_i$ sebuah zat terlarut mengikuti bentuk
yang sama, dengan difusivitas massa sebagai pengganti konduktivitas dan laju reaksi sebagai sumber.

Ketika perbedaan temperatur kecil, perubahan massa jenis bisa diabaikan di semua tempat kecuali
di suku gravitasi. Pendekatan Boussinesq menulis $\rho=\rho_0\big(1-\beta(T-T_0)\big)$ hanya di suku
gaya badan, dengan $\beta$ koefisien muai. Fluida yang lebih panas menjadi sedikit lebih ringan dan
naik. Pendekatan ini cukup untuk menjelaskan sel konveksi di panci yang dipanaskan, angin laut, dan
pancaran air hangat yang naik di dalam tangki.

\section{Batas dan awal}

Persamaan saja tidak menentukan sebuah aliran. Kita butuh syarat batas dan syarat awal. Di dinding
padat, fluida menempel dan kecepatannya nol, syarat yang disebut \istilah{no-slip}. Di saluran masuk,
kecepatan atau laju aliran ditentukan. Di saluran keluar, tekanan biasanya ditetapkan. Untuk
temperatur, dinding bisa diberi temperatur tetap, fluks panas tetap, atau dibuat adiabatik, yaitu
tanpa perpindahan panas sama sekali. Pilihan syarat batas sering lebih menentukan hasil simulasi
daripada skema numeriknya, dan seperti akan kita lihat, sebagian besar galat dalam membandingkan
simulasi dengan percobaan berasal dari sini.

\section{Mengapa simulasi aliran mahal}

Untuk aliran turbulen, pusaran terkecil yang membuang energi menjadi panas berukuran amat kecil
dibanding ukuran domain. Menyelesaikan semua skala itu secara langsung, yang disebut \istilah{direct
numerical simulation}, membutuhkan jumlah titik grid yang tumbuh kira-kira seperti $\mathrm{Re}^{9/4}$
dalam tiga dimensi, dan langkah waktu yang ikut mengecil \citep{pope2000}. Untuk aliran industri
dengan Re jutaan, itu jauh melampaui komputer mana pun. Simulasi praktis karena itu memodelkan efek
skala kecil alih-alih menyelesaikannya, dan pemodelan itulah yang membuka pintu bagi machine
learning: sebagai model untuk skala kecil, sebagai pengganti seluruh solver, atau sebagai cara
mempercepat solver yang ada. Bab-bab berikutnya membahas ketiganya.

\sumberbab{Bab ini mengikuti kuliah Numerical Methods in Fluid Mechanics bagian pengantar dan
persamaan atur (neraca massa, momentum, energi, spesies). Buku rujukan yang baik adalah
\citet{versteeg2007}, \citet{ferziger2020}, dan untuk turbulensi \citet{pope2000}.}
